\section{Token-level EEG 视觉特征增强}

在 A2 方案中，EEG 信息仅在 pooled embedding 层完成融合，所得的单条 512 维向量不再包含图像的空间分布信息。为了让 EEG 信号能够参与视觉语言模型原生的 token 级处理流程，本设计提出 VTF 方法，在 CLIP ViT 的 visual patch token 层级将 EEG 特征注入视觉表示。

\subsection{动机}

Pooled embedding fusion 的处理方式是将整幅图像压缩为一维向量后再与 EEG 特征合并，这种方法对图像分类场景基本够用，但在生成式任务中显露出不足：语言解码器天然以 token 序列作为输入，而非单个向量；而 token 序列本身保留了图像不同位置的空间视觉信息，有助于对图中多个实体进行详细描述。与此同时，EEG 对不同图像区域的神经响应可能存在差异——当受试者观看同时包含人物和物体的复杂画面时，各电极区域的激活模式往往各不相同。Token-level 融合利用注意力机制使各 visual patch token 能够有选择地关注不同的 EEG token，从而达成更细粒度的跨模态信息交换。

\subsection{CLIP ViT 视觉 token 提取}

CLIP ViT-B/32 将尺寸为 $224\times224$ 的输入图像划分为 $7\times7$ 个 patch（每个 patch 大小为 $32\times32$），经 ViT encoder 处理后生成 50 个视觉 token，其中包含 1 个 CLS token 和 49 个 patch token，各 token 维度均为 512，记作 $F_{\mathrm{vis}} \in \mathbb{R}^{B\times50\times512}$。CLS token 负责捕捉图像的全局语义，patch token 则编码各个局部区域的特征。

\subsection{EEG token 构造}

A2 encoder 输出的 pooled EEG embedding 是一条 512 维的单一向量。为了实现其与视觉 token 之间的 token 级交互，我们引入了一个轻量级的 EEG tokenizer：通过 4 组并行的可学习线性映射将 pooled embedding 扩展为 4 个 EEG token，即 $F_{\mathrm{eeg}} \in \mathbb{R}^{B\times4\times512}$。选取 4 个 token 是效率与信息多样性之间的权衡——token 数量太少会导致交互缺乏多样性，数量过多则可能使 EEG 信息过度主导视觉表示。

\subsection{视觉到 EEG 的交叉注意力融合}

VTF 的核心机制为 visual-to-EEG cross-attention，其整体结构如图~\ref{fig:vtf} 所示。在该设计中，视觉 token 充当 Query，EEG token 同时作为 Key 和 Value：

\begin{equation}
  F_{\mathrm{enhanced}} = F_{\mathrm{vis}} + \beta \cdot \mathrm{CrossAttention}(Q=F_{\mathrm{vis}}, K=F_{\mathrm{eeg}}, V=F_{\mathrm{eeg}})
\end{equation}

\placeholderfigure{VTF visual-to-EEG cross-attention 结构图}{vtf-cross-attention-placeholder}{4.2cm}

其中 CrossAttention 采用标准的多头交叉注意力机制（$h=8$）。$\beta$ 是一个置信度感知的门控系数，定义如下：

\begin{equation}
  \beta = \gamma \cdot \sigma\left(W_\beta [\mathrm{pool}(F_{\mathrm{vis}}); \mathrm{pool}(F_{\mathrm{eeg}})] + b_\beta\right)
\end{equation}

$\gamma$ 为可训练的缩放参数，初始值设为 0.1，目的是在训练早期抑制 EEG 注入量，防止梯度不稳定。该门控机制使模型能够依据视觉输入的质量动态调节 EEG 信息的注入强度。增强后的 $F_{\mathrm{enhanced}}$ 保持 $[B, 50, 512]$ 的形状不变，因此可以兼容任何以 CLIP ViT visual tokens 为输入的后续模型。

\subsection{分类验证}

为检验 VTF 的实际效果，我们首先在分类任务上进行评测：对 $F_{\mathrm{enhanced}}$ 做平均池化后，与 CLIP 文本原型进行零样本分类。结果显示，在多种退化条件下，real EEG 的表现均优于 vision-only 基线及 shuffled EEG（打乱脑电）/随机 EEG 对照，说明 token-level 的 EEG 融合是可行的。但 VTF 的整体分类准确率仍明显低于 A2 的 pooled fusion 方案，其原因是后者采用了更深的 gated fusion 网络并拥有更充分的跨模态交互。VTF 的主要价值在于其产生的 enhanced vision tokens 可以直接馈入生成式 VLM，这构成了下一章工作的基础。
